\documentclass[11pt]{article}
\usepackage{ich_guideline_style}
\addbibresource{references.bib}
\title{End-to-End Physical AI Oncology Clinical Trial Unification}
\author{Kevin Kawchak, ChemicalQDevice}
\date{12 March 2026}
\begin{document}

\begin{titlepage}
\centering
{\Large END-TO-END PHYSICAL AI ONCOLOGY CLINICAL TRIAL UNIFICATION\par}
\vspace{1.5cm}
{\Huge\bfseries Adaption: ICH Harmonised Guideline\par}
\vspace{0.5cm}
{\LARGE\bfseries Guideline for Good Physical AI Clinical Practice\par}
\vspace{0.3cm}
{\Large Modified E6(R3)\par}
\vspace{0.8cm}
{\large Draft release\par}
\vspace{0.2cm}
{\large Released on 12 March 2026\par}
\vspace{0.2cm}
{\large \href{https://doi.org/10.5281/zenodo.18973368}{10.5281/zenodo.18973368}\par}
\vspace{0.2cm}
{\large CEO Kevin Kawchak, ChemicalQDevice\par}
\vfill
{\small The original ICH E6(R3) document is protected by copyright and may be used, reproduced, incorporated into other works, adapted, modified, translated or distributed under a public license. ICH HARMONISED GUIDELINE: GUIDELINE FOR GOOD CLINICAL PRACTICE E6(R3) \href{https://database.ich.org/sites/default/files/ICH_E6\%28R3\%29_Step4_FinalGuideline_2025_0106.pdf}{(Source PDF)}. This current work is not endorsed or sponsored by ICH. The original PDF formatting was reconstructed into LaTeX, and adapted further using Claude Code Opus 4.6.\par}
\end{titlepage}

\pagenumbering{roman}
\tableofcontents
\clearpage
\section*{Prefatory Note}

This guidance has been adapted from the prior ICH E6(R3) regulation (adopted 06 January 2025) to address the end-to-end integration of physical AI, advanced robotics, and digital twin systems into oncology clinical trials. The prior ICH E6(R3) guideline established foundational principles for good clinical practice in traditional pharmaceutical trials. This current document extends those principles to encompass the specialized requirements of physical AI oncology clinical trial unification, including robotic surgical systems, collaborative robots (cobots), humanoid assistants, AI-driven monitoring, simulation-based validation, and federated multi-site trial coordination.

The Unification Standard Level (USL) scoring framework (v1.4.0 through v1.8.0) provides quantitative readiness assessments for robotic platforms across four dimensions: simulation switching, AI integration, cross-robot sharing, and clinical trial collaboration. USL scores range from 1.0 to 10.0 and are referenced throughout this guidance to establish minimum readiness thresholds for clinical deployment.

This guidance should be read in conjunction with the physical-ai-oncology-trials repository (v2.2.0), which provides production-ready configurations, validated pipelines, and integration guides for deploying robotic systems, digital twins, and embodied AI agents in oncology settings.

\section*{Document History}

\begin{description}[leftmargin=4cm,style=sameline]
\item[Prior Regulation] ICH E6(R3) Guideline for Good Clinical Practice, adopted 06 January 2025. The prior ICH E6(R3) regulation established the baseline framework for good clinical practice.
\item[Current Guidance] End-to-End Physical AI Oncology Clinical Trial Unification, released 12 March 2026. This guidance supersedes the prior ICH E6(R3) regulation for physical AI oncology trial contexts.
\item[Repository Reference] physical-ai-oncology-trials v2.2.0 (March 2026). Provides validated code, tools, and documentation supporting this guidance.
\end{description}

\clearpage
\pagenumbering{arabic}

\section*{Prior ICH E6(R3): Legal Notice}

This section preserves the legal notice from the prior ICH E6(R3) regulation. The original ICH E6(R3) document is protected by copyright and may be used, reproduced, incorporated into other works, adapted, modified, translated or distributed under a public license provided that ICH's copyright in the document is acknowledged at all times. In case of any adaptation, modification or translation of the document, reasonable steps must be taken to clearly label, demarcate or otherwise identify that changes were made to or based on the original document. Any impression that the adaptation, modification or translation of the original document is endorsed or sponsored by the ICH must be avoided. The document is provided ``as is'' without warranty of any kind. In no event shall the ICH or the authors of the original document be liable for any claim, damages or other liability arising from the use of the document. The above-mentioned permissions do not apply to content supplied by third parties. Therefore, for documents where the copyright vests in a third party, permission for reproduction must be obtained from this copyright holder.

\section{PRINCIPLES OF PHYSICAL AI CLINICAL PRACTICE}

The following principles establish the foundation for conducting oncology clinical trials that incorporate physical AI systems, advanced robotics, digital twins, and AI-driven decision support. These principles adapt and extend the prior ICH E6(R3) regulation to address the unique requirements of end-to-end physical AI unification in clinical settings.

\subsection{1.1 Foundational Principles}

\paragraph{1.1.1} Clinical trials involving physical AI systems should be conducted in accordance with the ethical principles that have their origin in the Declaration of Helsinki and that are consistent with good clinical practice (GCP) and the applicable regulatory requirements. Robotic interventions, AI-assisted diagnostics, and automated treatment delivery must preserve participant autonomy and welfare above all other considerations.

\paragraph{1.1.2} Before a physical AI clinical trial is initiated, foreseeable risks and inconveniences arising from robotic systems, AI algorithms, and automated processes should be weighed against the anticipated benefit for the individual trial participant and for oncology patients broadly. A trial should be initiated and continued only if the anticipated benefits justify the risks, including risks specific to robotic malfunction, AI prediction error, and cyber-physical system failure.

\paragraph{1.1.3} The rights, safety, and well-being of the trial participants are the most important considerations and should prevail over interests of science and society. This principle extends to interactions with all physical AI systems, including surgical robots (e.g., Intuitive da Vinci, Medtronic Hugo RAS, CMR Versius), cobots (e.g., Franka Emika Panda, Kinova Gen3, UFACTORY xArm 7), humanoid robots (e.g., Boston Dynamics Atlas, Tesla Optimus, Agility Digit), and all categories of therapeutic, diagnostic, assistive, and rehabilitative robotic systems.

\paragraph{1.1.4} Each individual involved in conducting a physical AI oncology trial should be qualified by education, training, and experience to perform their respective tasks, including operation and oversight of robotic systems, AI model interpretation, digital twin management, and simulation validation. Specialized training in physical AI systems must supplement traditional clinical trial training as defined by the prior ICH E6(R3) regulation.

\paragraph{1.1.5} Physical AI clinical trials should be scientifically sound and described in a clear, detailed protocol that incorporates specifications for all robotic systems, AI models, simulation frameworks, and digital twin configurations to be used. The protocol should reference applicable USL scores (v1.4.0 through v1.8.0) for each robotic platform and establish minimum readiness thresholds.

\paragraph{1.1.6} A trial should be conducted in compliance with the protocol that has received prior institutional review board (IRB) or independent ethics committee (IEC) approval, with specific review of physical AI components including robotic safety parameters, AI model validation evidence, human-robot interaction protocols, and cybersecurity measures.

\paragraph{1.1.7} The medical care given to and medical decisions made on behalf of participants should always be the responsibility of a qualified physician or, when appropriate, of a qualified dentist. Physical AI systems serve as tools under human oversight; autonomous robotic actions in clinical settings require continuous human monitoring and intervention capability as specified in the human oversight quality management framework (v1.2.0).

\paragraph{1.1.8} Freely given informed consent should be obtained from every trial participant prior to clinical trial participation and prior to any interaction with physical AI systems. Informed consent materials must clearly describe the nature, purpose, and risks of robotic interactions, AI-assisted decision making, and digital twin-based treatment planning. The consent process should include visual aids or demonstrations of robotic systems where practical.

\subsection{1.2 Physical AI System Classification for Clinical Trials}

For the purposes of this guidance, physical AI systems deployed in oncology clinical trials are classified according to the following categories, consistent with the patient instructions framework (v1.9.0 through v2.1.0) and USL scoring methodology:

\begin{description}[leftmargin=2.5cm,style=sameline]

\item[Surgical Robots] Teleoperated and semi-autonomous systems for tumor resection, biopsy, and surgical intervention. Evaluated platforms include da Vinci (dVRK, USL 7.1), Hugo RAS (USL 4.5), and Versius (USL 3.4). Minimum recommended USL for clinical deployment: 5.0 (Functional level).

\item[Cobots] Collaborative robots operating alongside clinical staff for sample handling, drug preparation, and laboratory automation. Evaluated platforms include Franka Emika Panda (USL 7.4), Kinova Gen3 (USL 5.7), and UFACTORY xArm 7 (USL 3.4). These systems must comply with ISO/TS 15066 collaborative workspace safety requirements.

\item[Humanoid Robots] Full-body robotic systems for patient assistance, logistics, and pediatric companionship. Evaluated platforms include Boston Dynamics Atlas (USL 5.8), Agility Digit (USL 4.2), and Tesla Optimus (USL 3.6). Clinical deployment requires additional human-robot interaction safety validation.

\item[Therapeutic Robots] Radiotherapy patient-positioning systems and motion-tracking robots for precise treatment delivery. These systems integrate with linear accelerator control and require real-time position verification within sub-millimeter tolerances.

\item[Diagnostic Robots] Robotic needle-placement systems, imaging assistant robots, and steerable needle systems for biopsy and tissue sampling. Integration with DICOM imaging pipelines and real-time ultrasound guidance is required.

\item[Assistive Robots] Social companion robots for pediatric oncology patients and logistical support systems. These systems prioritize psychological comfort and require age-appropriate interaction protocols.

\item[Rehabilitative Robots] Exoskeletons and robotic gait trainers for post-surgical recovery and functional restoration. These systems must accommodate patient-specific biomechanical parameters.

\end{description}

\subsection{1.3 AI and Machine Learning Framework Requirements}

Physical AI oncology trials utilize multiple categories of artificial intelligence and machine learning, each with distinct validation requirements. The following frameworks, referenced in the physical-ai-oncology-trials repository (v2.2.0), establish the computational foundation:

\begin{description}[leftmargin=2.5cm,style=sameline]

\item[Generative AI] Vision-Language-Action (VLA) models such as NVIDIA GR00T N1.6 for humanoid robot control, diffusion policies for surgical trajectory generation, and synthetic data generation via NVIDIA Cosmos Predict 2.5 for rare oncology scenarios. World models including Cosmos Reason 2 provide physics-aware simulation and reasoning through dual-system architecture.

\item[Agentic AI] LLM-based surgical assistants with multimodal perception, multi-agent coordination via CrewAI (v1.6.1) and LangGraph (v1.1.0) for multi-site trial management, natural language robot programming via ROS 2 Jazzy/Kilted integration, and standardized tool communication via Model Context Protocol (MCP) under the Linux Foundation AAIF.

\item[Reinforcement Learning] Sim-to-real transfer with domain randomization for surgical autonomy, hierarchical RL for complex surgical procedures, multi-agent RL for cooperative surgical assistance, and GPU-accelerated training via NVIDIA Isaac Lab (v2.3.1) reducing training time from days to hours.

\item[Self-Supervised Learning] Contrastive learning and foundation model pretraining for medical imaging feature extraction, enabling transfer learning across oncology imaging modalities without extensive manual annotation.

\item[Supervised Learning] Segmentation, detection, and classification models for tumor identification, surgical instrument tracking, and treatment response assessment. Integration with MONAI for medical imaging AI provides validated preprocessing and model architectures.

\end{description}

\subsection{1.4 Simulation and Digital Twin Requirements}

\subsubsection{1.4.1 Simulation Frameworks}

All physical AI systems intended for clinical deployment must undergo validated simulation testing prior to any patient contact. The following simulation frameworks, with established unification bridges (v1.0.0), are recognized for pre-clinical validation:

\begin{description}[leftmargin=2.5cm,style=sameline]

\item[(a)] NVIDIA Isaac Lab (v2.3.1) and Isaac Sim (v5.0.0) for GPU-accelerated robot training and high-fidelity physics simulation, with the Newton Physics Engine (Beta) for GPU physics co-developed by NVIDIA, DeepMind, and Disney.

\item[(b)] MuJoCo (v3.4.0) and MuJoCo Warp (Beta) for precision physics simulation with GPU optimization, providing complementary validation to Isaac-based training.

\item[(c)] Gazebo Sim (Jetty, v10.0.0) for ROS 2 integrated simulation with native sensor modeling, supporting hardware-in-the-loop testing.

\item[(d)] PyBullet (v3.2.5) for rapid prototyping and initial feasibility assessment of robotic behaviors.

\end{description}

The validated integration path for surgical robot training follows: NVIDIA Isaac Lab (training) to MuJoCo (validation) to dVRK Hardware (deployment) to Clinical use. Cross-framework validation via the unification bridge must demonstrate behavioral consistency with trajectory deviations below 2mm and force discrepancies below 0.5N across simulation engines.

\subsubsection{1.4.2 Digital Twin Systems}

Patient-specific digital twins provide computational models for treatment planning, response prediction, and virtual trial design. The digital twins framework (v1.0.0 through v1.3.0) supports the following capabilities:

\begin{description}[leftmargin=2.5cm,style=sameline]

\item[(a)] Patient-specific tumor modeling using reaction-diffusion equations calibrated to longitudinal imaging data, enabling tumor growth prediction and treatment response simulation.

\item[(b)] Treatment simulation using pharmacokinetic/pharmacodynamic (PK/PD) models for chemotherapy response prediction, radiation dose optimization, and immunotherapy modeling.

\item[(c)] Real-time intraoperative guidance through sensor fusion of imaging data, force measurements, and robotic telemetry, synchronized with the digital twin model.

\item[(d)] Virtual trial cohort generation for statistical power analysis and adaptive trial design, enabling in-silico pre-screening of protocol modifications.

\item[(e)] Validation and verification (V\&V) of digital twin predictions against clinical outcomes, with documented accuracy metrics and confidence intervals.

\end{description}

\subsection{1.5 Unification Standard Level (USL) Framework}

The USL framework (v1.4.0 through v1.8.0) provides a standardized scoring methodology for evaluating robotic platform readiness for unified, multi-site oncology clinical trials. USL scores are computed across four equally weighted dimensions:

\begin{description}[leftmargin=2.5cm,style=sameline]

\item[Dimension 1] Simulation Switching: Ability to transfer validated behaviors between simulation frameworks (Isaac, MuJoCo, Gazebo, PyBullet) using standardized model conversion (URDF, MJCF, SDF, USD).

\item[Dimension 2] AI Integration: Compatibility with generative AI, agentic AI, reinforcement learning, and supervised learning pipelines for clinical task execution.

\item[Dimension 3] Cross-Robot Sharing: Capability to share trained policies, calibration data, and operational procedures across different robot platforms within the same clinical category.

\item[Dimension 4] Clinical Trial Collaboration: Readiness for multi-site deployment with federated learning support, standardized data collection, and regulatory compliance documentation.

\end{description}

USL scores map to ten levels: Level 1 (Minimal, 1.0), Level 2 (Elementary, 2.0), Level 3 (Basic, 3.0), Level 4 (Developing, 4.0), Level 5 (Functional, 5.0), Level 6 (Proficient, 6.0), Level 7 (Advanced, 7.0), Level 8 (Expert, 8.0), Level 9 (Leading, 9.0), and Level 10 (Reference, 10.0). Three bands categorize readiness: Foundational (Levels 1 through 4), Intermediate (Levels 5 through 6), and Advanced (Levels 7 through 10).

\section{INVESTIGATOR RESPONSIBILITIES IN PHYSICAL AI TRIALS}

The investigator is responsible for the conduct of the physical AI oncology clinical trial at the trial site. If the investigator delegates trial-related duties involving physical AI systems, the investigator retains ultimate responsibility for the safety and integrity of all robotic interactions, AI-assisted decisions, and digital twin-based treatment plans. This section adapts the prior ICH E6(R3) investigator responsibilities to incorporate physical AI system oversight.

\subsection{2.1 Investigator Qualifications and Agreements}

\paragraph{2.1.1} The investigator should be qualified by education, training, and experience to assume responsibility for the proper conduct of a physical AI oncology trial, including demonstrated competency with all robotic systems, AI tools, and simulation platforms specified in the protocol. The investigator should meet all qualifications specified by the applicable regulatory requirements and should provide evidence of training in physical AI clinical systems.

\paragraph{2.1.2} The investigator should be thoroughly familiar with the appropriate use of the physical AI systems as described in the protocol, the Investigator's Brochure, and the Physical AI System Documentation package, including the current USL assessment, simulation validation reports, and safety parameter specifications for each robotic platform.

\paragraph{2.1.3} The investigator should be aware of and should comply with GCP requirements, the applicable regulatory requirements, and all physical AI system safety standards, including ISO/TS 15066 for collaborative robots, IEC 62304 for medical device software lifecycle, and IEC 80601-2-77 for robotically assisted surgical equipment.

\paragraph{2.1.4} The investigator and the sponsor should sign the protocol or an alternative document to confirm agreement, including specific provisions for physical AI system oversight, robotic safety monitoring, and AI model version control. The agreement should document the investigator's authority to pause or terminate robotic operations if participant safety concerns arise.

\subsection{2.2 Adequate Resources for Physical AI Trials}

\paragraph{2.2.1} The investigator should be able to demonstrate, at a minimum, the ability to recruit the required number of suitable trial participants within the agreed recruitment period, with consideration for participant acceptance of physical AI interactions.

\paragraph{2.2.2} The investigator should have sufficient time to properly conduct and complete the trial within the agreed trial period, including time required for physical AI system setup, calibration, simulation validation, and post-procedure data collection.

\paragraph{2.2.3} The investigator should have available an adequate number of qualified staff and adequate facilities for the foreseen duration of the trial, including:

\begin{description}[leftmargin=2.5cm,style=sameline]

\item[(a)] Trained robotic system operators certified on each platform specified in the protocol, with documented competency assessments and refresher training schedules.

\item[(b)] AI systems engineers responsible for model deployment, version control, performance monitoring, and anomaly detection during clinical procedures.

\item[(c)] Clinical engineering staff for hardware maintenance, calibration verification, and emergency response to robotic system failures.

\item[(d)] Adequate physical infrastructure including robotic operating theaters, shielded computation facilities, and network infrastructure meeting cybersecurity requirements.

\item[(e)] Backup systems and contingency procedures for all physical AI components, ensuring that no single point of failure can compromise participant safety.

\end{description}

\paragraph{2.2.4} The investigator should ensure that all persons assisting with the trial, including those operating physical AI systems, are adequately informed about the protocol, the investigational product(s), the robotic systems, and their trial-related duties and functions. A delegation log must document all physical AI system access permissions and operator qualifications.

\subsection{2.3 Medical Care of Trial Participants}

\paragraph{2.3.1} A qualified physician (or dentist, when appropriate) who is an investigator or a sub-investigator for the trial should be responsible for all trial-related medical decisions involving physical AI systems, including robotic procedure planning, AI-assisted diagnostic interpretation, and digital twin-based treatment optimization.

\paragraph{2.3.2} During and following a trial participant's interaction with physical AI systems, the investigator should ensure that adequate medical care is provided, including monitoring for adverse events specific to robotic interactions such as mechanical injury, unexpected AI-driven treatment modifications, or digital twin prediction discrepancies. The investigator should ensure that the participant is informed when medical care is needed for symptoms that may be related to physical AI system interactions.

\paragraph{2.3.3} The investigator should inform the trial participant's primary physician about the participant's involvement in the physical AI clinical trial, including the types of robotic systems encountered, if the participant agrees to such disclosure.

\subsection{2.4 Communication with IRB/IEC}

\paragraph{2.4.1} Before initiating a physical AI oncology trial, the investigator should have documented approval from the IRB/IEC for the trial protocol, informed consent forms, participant recruitment materials, and all Physical AI System Documentation packages. The IRB/IEC review should include evaluation of robotic safety data, AI model validation evidence, simulation test results, and human-robot interaction risk assessments.

\paragraph{2.4.2} The investigator should provide the IRB/IEC with all documents subject to review, including updates to robotic system software, AI model versions, and digital twin calibration parameters that may affect participant safety or trial outcomes.

\paragraph{2.4.3} The investigator should not implement any deviation from or changes to the protocol related to physical AI systems without agreement by the sponsor and prior review and documented approval from the IRB/IEC, except where necessary to eliminate an immediate hazard to trial participants arising from robotic system malfunction or AI system failure.

\paragraph{2.4.4} The investigator should promptly report to the IRB/IEC all changes to physical AI system configurations, AI model updates, robotic safety incidents, and any events that affect the risk-benefit assessment of the trial.

\subsection{2.5 Compliance with Protocol}

\paragraph{2.5.1} The investigator should conduct the trial in compliance with the protocol agreed to by the sponsor and approved by the IRB/IEC, including all specifications for physical AI system operation, AI model parameters, simulation validation criteria, and robotic safety thresholds.

\paragraph{2.5.2} The investigator should not implement a deviation from or a change of the protocol without agreement of the sponsor and prior review and documented approval from the IRB/IEC of an amendment, except where necessary to eliminate an immediate hazard to trial participants or when the change involves only logistical or administrative aspects that do not affect participant safety, the scientific integrity of the trial, or the physical AI system configurations.

\paragraph{2.5.3} The investigator should document and explain any deviation from the approved protocol, particularly deviations involving physical AI system parameters such as robotic force limits, AI model confidence thresholds, or digital twin prediction boundaries. Deviations from robotic safety parameters require immediate documentation and sponsor notification.

\subsection{2.6 Investigational Product and Physical AI Systems}

\paragraph{2.6.1} Responsibility for investigational product(s) accountability at the trial site(s) rests with the investigator, extending to all physical AI systems, robotic consumables, and AI-specific computational resources deployed at the site.

\paragraph{2.6.2} Where allowed by applicable regulatory requirements, the investigator may assign some or all of the physical AI system management duties to an appropriately trained clinical engineering team, such as a robotic systems pharmacist-equivalent or a clinical AI operations team. The investigator retains oversight responsibility.

\paragraph{2.6.3} The investigator and designated staff should maintain records of the physical AI systems at the site, including robotic system serial numbers, software versions, AI model identifiers, calibration dates, maintenance logs, and any modifications or updates applied during the trial period.

\subsection{2.7 Randomisation and Blinding in Physical AI Trials}

\paragraph{2.7.1} The investigator should follow the trial's randomisation procedures, if any, and should ensure that the code is broken only in accordance with the protocol. Where physical AI systems are involved in treatment allocation or dose optimization, the AI-driven randomisation processes must be validated against statistical requirements and the blinding integrity must be maintained across all automated systems.

\paragraph{2.7.2} In trials where robotic system type or AI algorithm variant is a blinded variable, the investigator should ensure that the blinding is maintained through appropriate physical and digital segregation of system identifiers, with documented procedures for emergency unblinding when participant safety requires knowledge of the specific system configuration.

\subsection{2.8 Informed Consent for Physical AI Interactions}

Informed consent for physical AI oncology trials must address additional elements beyond those specified in the prior ICH E6(R3) regulation. The following requirements supplement existing informed consent principles:

\paragraph{2.8.1} The language used in the informed consent form and in any other documented information provided to the participant should be as non-technical as practical and should be understandable to the participant. Descriptions of robotic systems, AI algorithms, and digital twins should use plain language supported by visual aids, photographs, or video demonstrations where feasible.

\paragraph{2.8.2} Before informed consent may be obtained, the participant should be given adequate time to review all physical AI system information, including the opportunity to observe robotic systems in non-clinical demonstration settings where available.

\paragraph{2.8.3} Neither the investigator nor the investigator site staff should coerce or unduly influence a participant to participate in or to continue participation in a trial involving physical AI systems. Participants should be informed that declining robotic-assisted procedures does not affect their access to standard-of-care treatment.

\paragraph{2.8.4} None of the information provided to the participant during the informed consent process should contain any language that causes the participant to waive or to appear to waive any legal rights, or that releases or appears to release the investigator, the institution, the sponsor, the robotic system manufacturer, or the AI model developer from liability for negligence.

\paragraph{2.8.5} The informed consent form for physical AI oncology trials should include, in addition to standard elements from the prior ICH E6(R3) regulation, the following:

\begin{description}[leftmargin=2.5cm,style=sameline]

\item[(a)] A description of each physical AI system the participant may encounter, including the type of robot, its clinical function, and the level of autonomy involved (teleoperated, shared autonomy, or supervised autonomous).

\item[(b)] An explanation of how AI algorithms contribute to treatment decisions, including the role of digital twin predictions in treatment planning and the possibility that AI recommendations may differ from standard clinical judgment.

\item[(c)] A description of the data collected by physical AI systems, including imaging data, force and motion telemetry, physiological measurements, and any data shared with the digital twin model. The participant should understand how this data is stored, protected, and used.

\item[(d)] A description of the safety measures in place, including human override capabilities, robotic system emergency stops, and the process for managing AI system failures or unexpected behaviors during clinical procedures.

\item[(e)] Information about the participant's right to request human-only (non-robotic) alternatives for any procedure, where such alternatives exist and are clinically appropriate.

\item[(f)] A description of any risks specific to physical AI systems, including mechanical risks from robotic contact, risks from AI prediction errors, cybersecurity risks, and any known limitations of the digital twin model being used.

\end{description}

\subsection{2.9 Records and Reports}

\paragraph{2.9.1} The investigator should maintain adequate and accurate records of all physical AI system interactions, including robotic procedure logs, AI model inference records, digital twin simulation outputs, and sensor data captured during clinical procedures. Source records for physical AI trials include traditional clinical documentation supplemented by electronic logs from robotic systems, AI decision audit trails, and digital twin synchronization records.

\paragraph{2.9.2} The investigator should ensure the accuracy, completeness, legibility, and timeliness of the data reported to the sponsor, including physical AI system performance metrics, robotic safety events, and any discrepancies between digital twin predictions and observed clinical outcomes.

\paragraph{2.9.3} Data changes in physical AI system records should be traceable, should not obscure the original entry, and should be explained where necessary. This applies to robotic telemetry corrections, AI model version updates, and digital twin recalibrations. An audit trail should be maintained for all computerised physical AI records.

\subsection{2.10 Safety Reporting in Physical AI Trials}

\paragraph{2.10.1} The investigator should report to the sponsor all serious adverse events (SAEs) immediately, including those arising from physical AI system interactions. Physical AI-specific reportable events include:

\begin{description}[leftmargin=2.5cm,style=sameline]

\item[(a)] Mechanical injury from robotic system contact, including unexpected forces, unplanned instrument movements, or collision events.

\item[(b)] Adverse outcomes attributable to AI-driven treatment decisions that deviated from clinical standard of care.

\item[(c)] Digital twin prediction failures where the simulated outcome differed materially from the observed clinical outcome, potentially affecting treatment decisions.

\item[(d)] Cybersecurity incidents affecting physical AI systems during clinical procedures, including unauthorized access, data corruption, or system compromise.

\item[(e)] Physical AI system malfunctions that interrupted clinical procedures or required emergency intervention, even if no participant injury resulted.

\end{description}

\paragraph{2.10.2} For reported deaths related to or potentially related to physical AI system interactions, the investigator should supply the sponsor and the IRB/IEC with any additional requested information, including complete robotic system logs, AI decision records, and digital twin state data at the time of the event.

\subsection{2.11 Premature Termination or Suspension of Physical AI Trial Components}

If the trial or specific physical AI components are prematurely terminated or suspended for any reason, the investigator should promptly inform the trial participants and should ensure that appropriate therapy and follow-up are provided. The investigator should inform the IRB/IEC and provide a detailed explanation, including whether the termination was caused by physical AI system safety concerns, AI model performance issues, or regulatory actions.

\subsection{2.12 Investigator Oversight of Physical AI Systems}

\paragraph{2.12.1} The investigator is responsible for the oversight of all physical AI activities at the investigator site. This oversight extends beyond traditional clinical trial supervision to encompass robotic system operations, AI model performance monitoring, digital twin accuracy assessment, and simulation validation verification.

\paragraph{2.12.2} The investigator should maintain a current delegation log that identifies all qualified personnel authorized to operate physical AI systems, interpret AI model outputs, manage digital twin configurations, and access robotic system controls. Delegated individuals must have documented training and competency assessment for each specific system.

\paragraph{2.12.3} The investigator should implement a systematic approach to physical AI system oversight, including:

\begin{description}[leftmargin=2.5cm,style=sameline]

\item[(a)] Regular review of robotic system performance metrics, including positional accuracy, force control precision, and procedure completion rates compared to simulation predictions.

\item[(b)] Periodic verification that AI models in clinical use match the validated versions specified in the protocol, with documented change control for any model updates.

\item[(c)] Ongoing assessment of digital twin prediction accuracy against observed clinical outcomes, with thresholds for triggering recalibration or model updates.

\item[(d)] Review of cybersecurity logs and access control records for all physical AI systems, with immediate investigation of any anomalous activity.

\item[(e)] Maintenance of emergency response protocols for physical AI system failures, including power loss, network disruption, and robotic system mechanical failure.

\end{description}

\paragraph{2.12.4} The investigator should ensure that the trial participants are informed of any significant changes to the physical AI systems used in their care, including software updates, hardware modifications, or changes to AI model versions, to the extent that such changes may affect the participant's willingness to continue in the trial.

\section{SPONSOR RESPONSIBILITIES IN PHYSICAL AI TRIALS}

The sponsor is responsible for implementing and maintaining quality assurance and quality control systems for all physical AI components of the trial. This section adapts the prior ICH E6(R3) sponsor responsibilities to address end-to-end physical AI oncology trial unification, including robotic system procurement, AI model lifecycle management, simulation validation, digital twin deployment, and federated multi-site coordination.

\subsection{3.1 Quality Management for Physical AI Trials}

\subsubsection{3.1.1 Critical to Quality Factors}

The sponsor should identify those factors that are critical to quality in the physical AI oncology trial. Critical to quality factors in physical AI trials extend beyond traditional clinical trial parameters to include:

\begin{description}[leftmargin=2.5cm,style=sameline]

\item[(a)] Robotic system safety parameters including force limits, workspace boundaries, speed constraints, and emergency stop response times.

\item[(b)] AI model performance thresholds including prediction accuracy, confidence calibration, inference latency, and drift detection sensitivity.

\item[(c)] Digital twin fidelity metrics including prediction accuracy for tumor growth, treatment response, and surgical outcomes relative to observed clinical data.

\item[(d)] Simulation validation criteria including cross-framework behavioral consistency as measured by trajectory deviation, force discrepancy, and task completion metrics.

\item[(e)] Data integrity for physical AI records including robotic telemetry, sensor data, AI inference logs, and digital twin state histories.

\item[(f)] Cybersecurity posture for all connected physical AI systems, including network segmentation, encryption standards, and intrusion detection capabilities.

\end{description}

\subsubsection{3.1.2 Risk Assessment and Proportionality}

The sponsor should apply a risk-based approach to quality management that considers the probability and impact of physical AI system failures on participant safety and trial reliability. Risk assessment for physical AI trials should incorporate:

\begin{description}[leftmargin=2.5cm,style=sameline]

\item[(a)] Robotic system failure modes and effects analysis (FMEA) for each clinical procedure involving robotic interaction, with documented mitigation strategies for each identified failure mode.

\item[(b)] AI model risk stratification based on the clinical consequence of prediction errors, with higher scrutiny for models that directly influence treatment decisions versus those used for monitoring or administrative purposes.

\item[(c)] Digital twin validation risk assessment based on the degree to which treatment decisions depend on twin predictions, with independent clinical verification required for high-consequence decisions.

\item[(d)] Cybersecurity threat modeling for the complete physical AI infrastructure, including robotic controllers, AI inference servers, digital twin computation platforms, and data communication networks.

\end{description}

\subsection{3.2 Regulatory Submission for Physical AI Components}

\paragraph{3.2.1} The sponsor should submit all necessary documentation to the regulatory authority(ies) and obtain approval before initiating the trial. For physical AI oncology trials, regulatory submissions should include documentation from the regulatory compliance framework (v1.0.0), covering FDA AI/ML guidance (January 2025), IEC 62304 medical device software lifecycle documentation, and ICH-GCP E6(R3) compliance verification adapted for physical AI contexts.

\paragraph{3.2.2} The sponsor should utilize the regulatory submission tools including the FDA submission tracker for 510(k), De Novo, PMA, and Breakthrough device pathways; the IRB protocol manager with AI-specific protocol elements; the GCP compliance checker (v1.0.0) for E6(R3) compliance verification; and the regulatory intelligence tracker for multi-jurisdiction monitoring.

\paragraph{3.2.3} The sponsor should inform the investigator(s) about regulatory decisions that affect the use of physical AI systems in the trial, including any conditions, restrictions, or additional requirements imposed by regulatory authorities on specific robotic systems, AI models, or digital twin methodologies.

\subsection{3.3 Confirmation of Review by IRB/IEC}

\paragraph{3.3.1} The sponsor should obtain from the IRB/IEC documented confirmation that the IRB/IEC has reviewed and approved the trial protocol, informed consent materials, and all Physical AI System Documentation packages, including robotic safety validation reports, AI model performance data, and digital twin accuracy assessments.

\paragraph{3.3.2} The IRB/IEC review for physical AI oncology trials should evaluate, at a minimum:

\begin{description}[leftmargin=2.5cm,style=sameline]

\item[(a)] The safety profile of each robotic system, including mechanical safety data, software reliability evidence, and USL score justification.

\item[(b)] The validation evidence for each AI model, including training data provenance, performance on representative oncology populations, and documented limitations.

\item[(c)] The accuracy and clinical relevance of digital twin predictions, including validation against independent clinical datasets.

\item[(d)] The adequacy of human oversight mechanisms, including human-in-the-loop controls, override capabilities, and emergency procedures.

\item[(e)] The cybersecurity measures protecting participant data and physical AI system integrity.

\item[(f)] The informed consent process, including participant understanding of physical AI interactions.

\end{description}

\subsection{3.4 Information on Physical AI Systems}

The sponsor should ensure that comprehensive documentation is developed and maintained for all physical AI systems used in the trial. The Physical AI System Documentation package is analogous to the Investigator's Brochure for investigational products, adapted from the prior ICH E6(R3) regulation Appendix A framework:

\begin{description}[leftmargin=2.5cm,style=sameline]

\item[(a)] For each robotic system: manufacturer, model, serial number, software version, USL score (with dimension-by-dimension breakdown), safety parameters, calibration records, maintenance schedule, and known limitations.

\item[(b)] For each AI model: model architecture, training data description (including provenance and de-identification methods), validation performance metrics, intended clinical use, known failure modes, version history, and update procedures.

\item[(c)] For each digital twin: computational model description, calibration methodology, validation evidence, prediction accuracy metrics, computational requirements, and data input specifications.

\item[(d)] For the simulation validation: framework versions used, cross-framework consistency results, behavioral comparison metrics, and any discrepancies noted during validation testing.

\item[(e)] For federated learning components: aggregation algorithm, differential privacy parameters (epsilon and delta values), secure aggregation protocol, and multi-site data harmonization procedures.

\end{description}

\subsection{3.5 Trial Design Considerations for Physical AI Trials}

\paragraph{3.5.1} The sponsor should engage appropriate qualified individuals, including robotic systems engineers, AI/ML scientists, clinical informaticists, and digital twin modelers, in the design of the trial protocol and supporting documentation.

\paragraph{3.5.2} Trial design for physical AI oncology trials should consider the following factors beyond traditional clinical trial design as described in the prior ICH E6(R3) regulation:

\begin{description}[leftmargin=2.5cm,style=sameline]

\item[(a)] Robot-specific variability: Different robotic platforms may exhibit different performance characteristics even when executing the same clinical task. The protocol should specify acceptable performance ranges and methods for controlling robot-specific confounding.

\item[(b)] AI model versioning: The protocol should specify which AI model versions are authorized for use and define the process for evaluating and approving model updates during the trial. Uncontrolled model updates represent a significant source of bias and are not permitted.

\item[(c)] Digital twin personalization: The protocol should specify the patient data required for digital twin calibration, the minimum acceptable prediction accuracy, and the process for handling cases where the digital twin cannot be adequately calibrated.

\item[(d)] Simulation-to-reality gap: The protocol should acknowledge the inherent difference between simulated and real-world robotic performance, specify the maximum acceptable sim-to-real transfer gap, and document the domain randomization parameters used during simulation training.

\item[(e)] Multi-site standardization: For multi-site trials, the protocol should specify standardization procedures for robotic system calibration, AI model deployment, and data collection to ensure cross-site comparability. Federated learning protocols should be documented in accordance with the federation framework (v1.0.0).

\end{description}

\subsection{3.6 Trial Management, Data Handling, and Record Keeping}

\paragraph{3.6.1} The sponsor should utilise appropriately qualified individuals to supervise the overall conduct of the physical AI trial, to handle the data, and to verify the data.

\paragraph{3.6.2} The sponsor should employ qualified robotic systems personnel, AI/ML engineers, and clinical data scientists to oversee the technical aspects of physical AI system deployment, monitoring, and data management.

\paragraph{3.6.3} The sponsor may consider establishing an independent Physical AI Safety Monitoring Committee, analogous to the Independent Data Monitoring Committee (IDMC), with expertise in robotic surgery, AI safety, and digital twin validation, to provide ongoing safety oversight for the physical AI components of the trial.

\subsection{3.7 Investigator Selection for Physical AI Trials}

\paragraph{3.7.1} The sponsor should select investigators qualified by training, experience, and resources to conduct the physical AI oncology trial, including demonstrated competency with the specific robotic systems, AI tools, and simulation platforms specified in the protocol.

\paragraph{3.7.2} Before entering an agreement with an investigator to conduct a physical AI trial, the sponsor should assess the investigator site's physical AI readiness, including:

\begin{description}[leftmargin=2.5cm,style=sameline]

\item[(a)] Robotic system infrastructure: availability, condition, and calibration status of required robotic platforms.

\item[(b)] Computing infrastructure: GPU computation capacity for AI model inference, digital twin simulation, and data processing.

\item[(c)] Network infrastructure: bandwidth, latency, and cybersecurity posture sufficient for physical AI system communications and federated learning participation.

\item[(d)] Staff qualifications: documented training and competency for all personnel who will operate or oversee physical AI systems.

\item[(e)] Emergency preparedness: documented procedures for managing physical AI system failures, including backup systems and manual fallback procedures.

\end{description}

\subsection{3.8 Allocation of Responsibilities}

\paragraph{3.8.1} The sponsor may transfer or delegate any trial-related duties and functions related to physical AI systems to a service provider, but the ultimate responsibility for the quality and integrity of the trial data remains with the sponsor, consistent with the prior ICH E6(R3) regulation.

\paragraph{3.8.2} Service providers performing physical AI system functions (e.g., robotic system maintenance, AI model training, digital twin development, simulation validation) should be selected, assessed, and overseen in accordance with the sponsor's quality management framework. Documentation of service provider qualifications should include demonstrated expertise in the specific physical AI technologies used in the trial.

\paragraph{3.8.3} Any transfer of physical AI system responsibilities to a service provider should be documented in writing, with clear specification of the transferred duties, quality standards, reporting requirements, and the sponsor's oversight mechanisms.

\subsection{3.9 Compensation to Participants and Investigators}

\paragraph{3.9.1} If required by the applicable regulatory requirements, the sponsor should provide insurance or should indemnify the investigator and institution against claims arising from the trial, including claims arising from physical AI system interactions, robotic malfunction, AI-driven treatment decisions, or digital twin prediction errors.

\paragraph{3.9.2} The sponsor's policies and procedures should address the costs of treatment of trial participants in the event of injuries related to physical AI systems, including injuries arising from robotic system contact, AI-guided treatment modifications, or medical device cybersecurity incidents.

\subsection{3.10 Monitoring Physical AI Trial Conduct}

\subsubsection{3.10.1 Purpose of Monitoring}

The purpose of trial monitoring in physical AI oncology trials is to verify that the rights and well-being of trial participants are protected, that the reported trial data are accurate and complete, and that all physical AI systems are operating within their validated parameters. Monitoring extends the prior ICH E6(R3) monitoring framework to encompass robotic system performance, AI model behavior, digital twin accuracy, and simulation validation compliance.

\subsubsection{3.10.2 Selection and Qualifications of Monitors}

Monitors assigned to physical AI oncology trials should have the scientific and clinical knowledge, as well as the physical AI technical expertise, needed to provide adequate monitoring of the trial. This includes understanding of robotic system operations, AI model lifecycle management, digital twin methodology, and relevant cybersecurity principles.

\subsubsection{3.10.3 Monitoring Plan for Physical AI Components}

The sponsor should develop a monitoring plan that addresses the unique requirements of physical AI trial components. The monitoring plan should include:

\begin{description}[leftmargin=2.5cm,style=sameline]

\item[(a)] Robotic system monitoring: scheduled verification of calibration, safety parameter compliance, and maintenance adherence.

\item[(b)] AI model monitoring: continuous or periodic assessment of model performance against validation benchmarks, including drift detection and performance degradation alerts.

\item[(c)] Digital twin monitoring: comparison of twin predictions against accumulating clinical outcomes, with pre-specified recalibration triggers.

\item[(d)] Centralised monitoring: use of data analytics tools from the physical-ai-oncology-trials repository (v2.2.0) for cross-site comparison of physical AI system performance, including the trial site monitor, deployment readiness checker, and simulation job runner.

\item[(e)] Cybersecurity monitoring: ongoing assessment of physical AI system security posture, including intrusion detection, access control verification, and incident response readiness.

\end{description}

\subsubsection{3.10.4 Monitoring Activities}

On-site and remote monitoring activities for physical AI trials should include, in addition to standard clinical trial monitoring from the prior ICH E6(R3) regulation:

\begin{description}[leftmargin=2.5cm,style=sameline]

\item[(a)] Verification that the correct robotic systems, with the validated software versions, are being used for each clinical procedure as specified in the protocol.

\item[(b)] Review of robotic system logs to confirm that safety parameters were maintained throughout each procedure, including force limits, workspace boundaries, and speed constraints.

\item[(c)] Verification that AI model versions in clinical use match the protocol-specified versions, with review of any model updates and their approval documentation.

\item[(d)] Assessment of digital twin prediction accuracy by comparing recent predictions against observed outcomes.

\item[(e)] Review of cybersecurity incident logs and verification that access controls are properly maintained for all physical AI systems.

\item[(f)] Verification that investigator site staff maintain current qualifications for physical AI system operation.

\end{description}

\subsubsection{3.10.5 Monitoring Report}

Reports of monitoring activities for physical AI trials should include a summary of what was reviewed for both traditional clinical and physical AI components, a description of significant findings, conclusions and actions required to resolve them, and follow-up on their resolution. Physical AI-specific monitoring report elements should include robotic system performance summaries, AI model status reports, digital twin accuracy assessments, and cybersecurity posture evaluations.

\subsection{3.11 Noncompliance in Physical AI Trials}

\paragraph{3.11.1} Noncompliance with the protocol, SOPs, GCP, the physical AI system specifications, or applicable regulatory requirements by an investigator or by the sponsor's staff should lead to appropriate and proportionate action by the sponsor to secure compliance. Physical AI-specific noncompliance events include operation of robotic systems outside validated parameters, use of unauthorized AI model versions, failure to maintain digital twin calibration, and cybersecurity standard violations.

\paragraph{3.11.2} If noncompliance that significantly affects or has the potential to significantly affect the rights, safety, or well-being of trial participants or the reliability of trial results is discovered, including noncompliance related to physical AI system specifications, the sponsor should perform a root cause analysis, implement appropriate corrective and preventive actions, and confirm their adequacy. Where the sponsor identifies serious noncompliance, the sponsor should notify the regulatory authority and IRB/IEC in accordance with applicable regulatory requirements.

\subsection{3.12 Safety Assessment and Reporting for Physical AI Systems}

The sponsor is responsible for the ongoing safety evaluation of all physical AI systems used in the trial, in addition to the safety evaluation of the investigational product(s) as specified in the prior ICH E6(R3) regulation.

\subsubsection{3.12.1 Sponsor Review of Physical AI Safety Information}

The sponsor should aggregate and review in a timely manner all safety information related to physical AI system operations, including:

\begin{description}[leftmargin=2.5cm,style=sameline]

\item[(a)] Robotic system incident reports from all trial sites, aggregated to identify trends, common failure modes, and potential systemic issues across the robotic fleet.

\item[(b)] AI model performance data across sites, with attention to site-specific performance variations that may indicate calibration issues, data quality problems, or population-specific model limitations.

\item[(c)] Digital twin prediction accuracy trends, with analysis of whether prediction errors are systematic or random and whether recalibration is needed.

\item[(d)] Cybersecurity threat intelligence relevant to the physical AI systems in use, including newly discovered vulnerabilities and recommended patches.

\end{description}

\subsubsection{3.12.2 Physical AI Safety Reporting}

\begin{description}[leftmargin=2.5cm,style=sameline]

\item[(a)] The sponsor should submit to the regulatory authority(ies) safety updates and periodic reports related to physical AI systems, including aggregate robotic safety data, AI model performance trends, and digital twin accuracy assessments.

\item[(b)] The sponsor should expedite reporting of physical AI events that constitute suspected serious safety concerns, including robotic system malfunctions resulting in participant injury, AI-driven treatment decisions resulting in adverse outcomes, and cybersecurity breaches affecting clinical systems.

\item[(c)] The sponsor should communicate relevant physical AI safety information to all investigators in a timely manner, including safety alerts for specific robotic system models, AI model recalls, and cybersecurity advisories.

\end{description}

\subsubsection{3.12.3 Managing Immediate Physical AI Hazards}

The sponsor should take prompt action to address immediate hazards arising from physical AI systems, including ordering the suspension of specific robotic system operations, recalling specific AI model versions, or isolating affected digital twin systems. The sponsor should determine the root causes and take appropriate remedial actions. The information on the immediate hazard should be communicated to the IRB/IEC, regulatory authorities, and investigators in a timely manner.

\subsection{3.13 Investigational Product and Physical AI System Supply}

\subsubsection{3.13.1 Information on Physical AI Systems}

The sponsor should ensure that a Physical AI System Documentation package is developed and updated as significant new information becomes available, analogous to the Investigator's Brochure requirements from the prior ICH E6(R3) regulation Appendix A.

\subsubsection{3.13.2 Physical AI System Preparation and Deployment}

\begin{description}[leftmargin=2.5cm,style=sameline]

\item[(a)] The sponsor should ensure that all physical AI systems are manufactured, tested, and validated in accordance with applicable device standards, including ISO 13485 (quality management for medical devices), IEC 62304 (medical device software lifecycle), and IEC 80601-2-77 (robotically assisted surgical equipment).

\item[(b)] The sponsor should determine acceptable operating parameters for each physical AI system, including environmental conditions, calibration requirements, and maintenance schedules. The sponsor should inform all involved parties of these requirements.

\item[(c)] The sponsor should ensure that software updates for physical AI systems are managed through a documented change control process, with validation evidence for each update before deployment to clinical sites.

\end{description}

\subsubsection{3.13.3 Supplying and Maintaining Physical AI Systems}

\begin{description}[leftmargin=2.5cm,style=sameline]

\item[(a)] The sponsor is responsible for ensuring that investigator sites have access to properly maintained physical AI systems throughout the trial. For robotic systems, this includes timely provision of replacement parts, software updates, and technical support.

\item[(b)] The sponsor should maintain records documenting the deployment, configuration, maintenance, and decommissioning of physical AI systems at each trial site, including system serial numbers, software version histories, and calibration records.

\item[(c)] The sponsor should establish processes for managing the end-of-trial disposition of physical AI systems, including data extraction, secure erasure of participant data from robotic controllers and AI inference engines, and system return or repurposing procedures.

\end{description}

\subsection{3.14 Data and Records for Physical AI Trials}

\subsubsection{3.14.1 Data Handling}

\begin{description}[leftmargin=2.5cm,style=sameline]

\item[(a)] The sponsor should ensure the integrity and confidentiality of all data generated by physical AI systems, including robotic telemetry, AI inference records, digital twin state data, and sensor measurements. The privacy framework (v1.0.0 through v1.3.0) provides HIPAA-compliant tools for PHI/PII detection, de-identification, access control, breach response, and data use agreement generation.

\item[(b)] The sponsor should apply quality control at relevant stages of physical AI data handling to ensure that the data are of sufficient quality to generate reliable results. Quality control should focus on data of higher criticality, including safety-relevant robotic measurements, AI model decision records, and digital twin predictions that influenced treatment decisions.

\item[(c)] The sponsor should pre-specify in the protocol what physical AI data will be collected and the method of collection, including the sampling rates for robotic telemetry, the granularity of AI decision logging, and the frequency of digital twin state captures.

\item[(d)] The sponsor should ensure that physical AI data acquisition systems are validated and ready for use prior to their deployment in the trial, including robotic data loggers, AI inference monitoring tools, and digital twin synchronization systems.

\item[(e)] The sponsor should ensure that documented processes are implemented to maintain data integrity for the full physical AI data lifecycle, from capture during clinical procedures through analysis and long-term archival.

\end{description}

\subsubsection{3.14.2 Statistical Programming and Data Analysis}

This section should be read in conjunction with the prior ICH E6(R3) regulation section on statistical programming and ICH E9 Statistical Principles for Clinical Trials, with additional considerations for physical AI data:

\begin{description}[leftmargin=2.5cm,style=sameline]

\item[(a)] The statistical analysis plan should address how physical AI system data will be incorporated into the analysis, including robotic performance covariates, AI model prediction accuracy as a secondary endpoint, and digital twin calibration quality as a data quality indicator.

\item[(b)] The sponsor should ensure appropriate quality control of statistical programming for physical AI data, including validation of data extraction pipelines from robotic systems, verification of AI decision log parsing algorithms, and reconciliation of digital twin prediction records with clinical outcome data.

\item[(c)] The sponsor should ensure traceability of data transformations involving physical AI data, from raw sensor measurements through derived clinical variables.

\end{description}

\subsubsection{3.14.3 Record Keeping and Retention}

\begin{description}[leftmargin=2.5cm,style=sameline]

\item[(a)] The sponsor should retain all physical AI-specific essential records pertaining to the trial, including robotic system logs, AI model artifacts and version histories, digital twin models and calibration data, simulation validation reports, and cybersecurity incident records, in accordance with the applicable regulatory requirements.

\item[(b)] The sponsor should inform investigators of the retention requirements for physical AI system records and should ensure that investigators can maintain access to physical AI data for the required retention period, even if the physical AI systems are decommissioned or returned.

\end{description}

\subsubsection{3.14.4 Record Access}

\begin{description}[leftmargin=2.5cm,style=sameline]

\item[(a)] The sponsor should ensure that direct access to physical AI system records is available for trial-related monitoring, audits, regulatory inspections, and IRB/IEC review. Physical AI records should be provided in accessible formats with appropriate documentation of data schemas and metadata.

\item[(b)] The sponsor should ensure that trial participants have consented to direct access to physical AI system records that contain data from their clinical interactions.

\end{description}

\subsection{3.15 Clinical Trial Reports for Physical AI Trials}

\subsubsection{3.15.1 Premature Termination or Suspension}

If a trial or its physical AI components are prematurely terminated or suspended, the sponsor should promptly inform the investigators, institutions, and regulatory authorities of the termination or suspension and the reasons, including whether the termination was related to physical AI system safety concerns. The sponsor should provide investigators with information about alternative therapy and follow-up for participants.

\subsubsection{3.15.2 Clinical Trial Reports}

\begin{description}[leftmargin=2.5cm,style=sameline]

\item[(a)] Whether the trial is completed or prematurely terminated, the sponsor should ensure that the clinical trial reports include comprehensive documentation of physical AI system performance, including aggregate robotic safety data, AI model performance summaries, digital twin prediction accuracy analyses, and simulation validation results.

\item[(b)] The clinical trial report should document any physical AI system-related protocol deviations, safety events, or unexpected findings that may have affected trial outcomes.

\item[(c)] Once the trial has been completed and relevant analyses finalized, the sponsor should make physical AI-specific trial results publicly available, including robotic system performance benchmarks and AI model validation results, in accordance with applicable regulatory requirements.

\end{description}

\section{DATA GOVERNANCE FOR PHYSICAL AI TRIALS -- INVESTIGATOR AND SPONSOR}

This section provides guidance to the responsible parties (investigators and sponsors) on appropriate management of data integrity, traceability, and security for physical AI oncology trials, thereby allowing the accurate reporting, verification, and interpretation of trial-related information. This section adapts the prior ICH E6(R3) data governance framework (Section 4) to address the specialized data requirements of robotic telemetry, AI model inference records, digital twin state data, and federated multi-site data sharing.

The quality and amount of information generated in a physical AI clinical trial should be sufficient to address trial objectives, provide confidence in the trial's results, and support good decision making. Physical AI data sources include traditional clinical data, robotic system telemetry, AI decision logs, digital twin models, simulation records, and sensor measurements.

The systems and processes that help ensure this quality should be designed and implemented in a way that is proportionate to the risks to participants and the reliability of trial results, with additional considerations for the volume, velocity, and complexity of physical AI data streams.

The following key processes should address the full data lifecycle for physical AI trials, with a focus on the criticality of the data, and should be implemented proportionately and documented appropriately:

\begin{description}[leftmargin=2.5cm,style=sameline]

\item[(a)] Processes to ensure the protection of the confidentiality of trial participants' data collected by physical AI systems, including robotic imaging, force telemetry, physiological measurements, and AI-processed biomarkers. The privacy framework (v1.0.0) provides HIPAA-compliant tools for PHI/PII detection, de-identification using Safe Harbor and Expert Determination methods, role-based access control with 21 CFR Part 11 audit trails, and automated breach response protocols.

\item[(b)] Processes for managing computerised systems, including robotic controllers, AI inference engines, digital twin platforms, and simulation environments, to ensure that they are fit for purpose and used appropriately.

\item[(c)] Processes to safeguard essential elements of the clinical trial, such as randomisation, dose adjustments, blinding, and physical AI system parameters that affect treatment delivery.

\item[(d)] Processes to support key decision making, including data finalisation prior to analysis, unblinding, allocation to analysis data sets, changes in trial design, and the activities of the IDMC or Physical AI Safety Monitoring Committee.

\end{description}

\subsection{4.1 Safeguard Blinding in Physical AI Data Governance}

\paragraph{4.1.1} Maintaining the integrity of the blinding is important in the design of physical AI systems, management of user accounts for robotic controllers and AI platforms, delegation of responsibilities with respect to physical AI data handling, and provision of data access at sites. Physical AI systems may inadvertently reveal treatment assignments through differences in robotic behavior, AI model outputs, or digital twin configurations. The sponsor should assess and mitigate these risks.

\paragraph{4.1.2} Roles, responsibilities, and procedures for access to unblinded information in physical AI systems should be defined and documented by all relevant parties according to the protocol. In blinded trials, robotic system operators and AI engineers who interact with investigator site staff should not have access to unblinding information except when justified by the trial design.

\paragraph{4.1.3} Suitable mitigation strategies should be implemented to reduce the risk of inadvertent unblinding through physical AI system data, including masking of treatment-specific robotic parameters in system displays, anonymisation of AI model variant identifiers, and restriction of digital twin configuration details.

\paragraph{4.1.4} The potential for unblinding through physical AI data should be part of the risk assessment of a blinded trial. Any planned or unplanned unblinding, including inadvertent exposure of physical AI system configurations, should be documented and assessed for its impact on the trial results.

\subsection{4.2 Data Lifecycle Elements for Physical AI Trials}

Procedures should be in place to cover the full data lifecycle for all physical AI data, from initial capture during clinical procedures through analysis and long-term retention.

\subsubsection{4.2.1 Data Capture}

\begin{description}[leftmargin=2.5cm,style=sameline]

\item[(a)] Physical AI data capture encompasses multiple modalities: robotic joint positions, velocities, and torques captured at control frequencies (typically 50 to 1000 Hz); force and torque measurements from instrument-tissue interactions; imaging data from endoscopic, ultrasound, or CT guidance systems; AI model inputs, outputs, and confidence scores; digital twin state vectors and prediction outputs; and environmental sensor data including temperature, humidity, and electromagnetic interference levels.

\item[(b)] Acquired data from physical AI systems should be accompanied by relevant metadata, including timestamps synchronized to coordinated universal time (UTC), system identifiers, software versions, calibration status, and operator identifiers.

\item[(c)] At the point of data capture, automated validation checks should be implemented to detect sensor malfunctions, data transmission errors, and physical AI system anomalies. These checks should be proportionate to the criticality of the data being captured.

\end{description}

\subsubsection{4.2.2 Relevant Metadata, Including Audit Trails}

The approach used by the responsible party for implementing, evaluating, accessing, managing, and reviewing relevant metadata associated with physical AI data of higher criticality should include:

\begin{description}[leftmargin=2.5cm,style=sameline]

\item[(a)] Evaluating the physical AI system for the types and content of metadata available to ensure that:

\item[(i)] Robotic control systems maintain logs of operator account creation, permission changes, and access events.

\end{description}

(ii) AI inference engines record model version, input data characteristics, output predictions, confidence scores, and any human override actions.

(iii) Digital twin platforms maintain complete state histories, parameter changes, calibration events, and prediction records.

(iv) All physical AI systems record workflow actions, configuration changes, and safety-relevant events in addition to direct data entries.

\begin{description}[leftmargin=2.5cm,style=sameline]

\item[(b)] Ensuring that audit trails for physical AI systems are not disabled. Robotic system logs, AI inference records, and digital twin state histories should be immutable and tamper-evident. Audit trail modifications should occur only in documented exceptional circumstances.

\item[(c)] Ensuring that physical AI audit trails are interpretable by qualified reviewers, with appropriate documentation of data formats, units of measurement, coordinate systems, and reference frames.

\item[(d)] Ensuring that automatic capture of date and time for all physical AI data entries and transfers is unambiguous, using coordinated universal time (UTC) with sufficient precision for the data type (millisecond precision for robotic telemetry, second precision for administrative records).

\item[(e)] Determining which physical AI metadata require review and long-term retention, based on the criticality of the data and its role in supporting trial conclusions.

\end{description}

\subsubsection{4.2.3 Review of Physical AI Data and Metadata}

Procedures for review of physical AI trial data, audit trails, and other relevant metadata should be in place. Review should be a planned activity, and the extent and nature should be risk-based, adapted to the individual trial, and adjusted based on experience during the trial. Physical AI data review should include periodic assessment of robotic system performance trends, AI model drift detection, and digital twin prediction accuracy tracking.

\subsubsection{4.2.4 Data Corrections}

There should be processes to correct errors in physical AI data that could impact the reliability of the trial results. Corrections should be attributed to the person or computerised system making the correction, justified, and supported by source records. For robotic telemetry data, corrections should document the original sensor reading, the corrected value, and the technical justification for the correction (e.g., sensor calibration offset, known measurement artifact).

\subsubsection{4.2.5 Data Transfer, Exchange, and Migration}

Validated processes should be in place to ensure that physical AI data, including relevant metadata, transferred between systems retains its integrity and preserves its confidentiality. This includes:

\begin{description}[leftmargin=2.5cm,style=sameline]

\item[(a)] Transfer of robotic telemetry from surgical systems to data management platforms.

\item[(b)] Transfer of AI model inference records from edge computing devices to centralised analysis systems.

\item[(c)] Synchronisation of digital twin state data between clinical sites and centralised twin computation platforms.

\item[(d)] Federated data exchange between trial sites using the federation framework (v1.0.0), including differential privacy mechanisms, secure aggregation protocols, and data harmonization procedures for DICOM, FHIR, ICD-10 to SNOMED CT, and LOINC tumor marker standardisation.

\end{description}

The data exchange and transfer process should be documented to ensure traceability, and data reconciliation should be implemented to avoid data loss and unintended modifications.

\subsubsection{4.2.6 Finalisation of Data Sets Prior to Analysis}

\begin{description}[leftmargin=2.5cm,style=sameline]

\item[(a)] Data of sufficient quality for interim and final analysis should be defined and are achieved by implementing timely and reliable processes for physical AI data capture, verification, validation, review, and rectification of errors. Physical AI data quality assessment should include verification of sensor calibration status, AI model version consistency, and digital twin synchronisation integrity.

\item[(b)] Activities undertaken to finalise the physical AI data sets prior to analysis should be confirmed and documented. These activities may include reconciliation of robotic telemetry with clinical outcome records, verification of AI model version logs, reconciliation of digital twin predictions with observed outcomes, and resolution of data quality flags raised during monitoring.

\item[(c)] Data extraction from physical AI systems and determination of analysis data sets should take place in accordance with the planned statistical analysis and should be documented, including the data extraction queries, filtering criteria, and any transformations applied.

\end{description}

\subsubsection{4.2.7 Retention and Access}

Physical AI trial data and relevant metadata should be archived in a way that allows for their retrieval and readability and should be protected from unauthorised access and alterations throughout the retention period. For physical AI data, retention should address the preservation of data format specifications, coordinate system definitions, and software tools needed to interpret archived robotic telemetry, AI model records, and digital twin state data.

\subsubsection{4.2.8 Destruction}

Physical AI trial data and metadata may be permanently destroyed when no longer required as determined by applicable regulatory requirements. Destruction of physical AI data should follow secure erasure procedures, particularly for data stored on robotic controllers, edge computing devices, and AI inference platforms that may retain cached data.

\subsection{4.3 Computerised Systems for Physical AI Trials}

As described in sections 2 and 3, the responsibilities of the sponsor, investigator, and other parties with respect to computerised systems used in physical AI clinical trials should be clear and documented. Physical AI trials involve a broader range of computerised systems than traditional clinical trials, including robotic control systems, AI inference engines, digital twin computation platforms, simulation environments, and federated learning infrastructure.

The responsible party should ensure that those developing or configuring physical AI systems for clinical trials are aware of the intended clinical purpose and the regulatory requirements that apply.

It is recommended that representatives of intended participant populations and healthcare professionals are involved in the design and evaluation of physical AI systems intended for direct patient interaction, to ensure that these systems are suitable for use by the intended patient population.

\subsubsection{4.3.1 Procedures for the Use of Physical AI Systems}

Documented procedures should be in place to ensure the appropriate use of physical AI systems in clinical trials. These procedures should cover:

\begin{description}[leftmargin=2.5cm,style=sameline]

\item[(a)] Robotic system startup, calibration verification, and operational readiness checks before each clinical use.

\item[(b)] AI model deployment, version verification, and performance validation before clinical inference.

\item[(c)] Digital twin initialisation, patient data loading, calibration, and prediction generation.

\item[(d)] Simulation environment setup, parameter configuration, and cross-framework validation execution.

\item[(e)] Federated learning round initiation, model aggregation, and privacy budget management.

\end{description}

\subsubsection{4.3.2 Training}

The responsible party should ensure that those using physical AI systems are appropriately trained in their use. Training should cover system-specific operation, safety procedures, emergency response, data handling, and the clinical context of the trial. Training records should be maintained and should demonstrate competency, not merely attendance. Refresher training should be provided when system updates alter the user interface or operational procedures.

\subsubsection{4.3.3 Security}

\begin{description}[leftmargin=2.5cm,style=sameline]

\item[(a)] The security of physical AI trial data and records should be managed throughout the data lifecycle, with particular attention to the expanded attack surface created by networked robotic systems, edge computing devices, and cloud-based AI inference services.

\item[(b)] The responsible party should ensure that security controls are implemented and maintained for all physical AI systems. These controls should include network segmentation between robotic control networks and clinical data networks, encryption of data in transit and at rest, multi-factor authentication for system access, and continuous monitoring for anomalous activity. Cybersecurity measures should align with the NIST Cybersecurity Framework and FDA premarket cybersecurity guidance.

\item[(c)] The responsible party should maintain adequate backup of all physical AI data, including robotic configuration data, AI model artifacts, and digital twin state histories.

\item[(d)] Procedures should cover disaster recovery for physical AI systems, including documented procedures for restoring robotic systems, AI inference capabilities, and digital twin platforms following system failure. Recovery time objectives should be defined based on clinical impact.

\end{description}

\subsubsection{4.3.4 Validation}

\begin{description}[leftmargin=2.5cm,style=sameline]

\item[(a)] The responsible party is responsible for the validation status of physical AI systems throughout their lifecycle. The approach to validation should be based on a risk assessment that considers the intended clinical use, the criticality of the data generated, and the potential of the system to affect the well-being, rights, and safety of trial participants and the reliability of trial results.

\item[(b)] Validation of physical AI systems should demonstrate that each system conforms to established requirements for completeness, accuracy, and reliability. For robotic systems, this includes positional accuracy, force control precision, and safety response time. For AI models, this includes prediction accuracy, calibration, fairness across demographic subgroups, and robustness to input perturbations. For digital twins, this includes prediction accuracy relative to observed outcomes.

\item[(c)] Physical AI systems should be validated prior to clinical use. Subsequent changes, including software updates, hardware replacements, and AI model retraining, should be validated based on risk and should consider both previously collected and new data in line with change control procedures.

\item[(d)] Periodic review should be conducted to ensure that physical AI systems remain in a validated state, including monitoring for performance degradation, hardware wear, and AI model drift.

\item[(e)] Cross-framework validation should verify that policies trained in one simulation environment (e.g., NVIDIA Isaac Lab v2.3.1) produce equivalent behaviors when transferred to a validation environment (e.g., MuJoCo v3.4.0), with documented tolerance thresholds.

\item[(f)] Validation documentation for physical AI systems should be maintained and retained, including test plans, test results, acceptance criteria, and any deviations or waivers.

\end{description}

\subsubsection{4.3.5 System Release}

Physical AI systems, including robotic configurations, AI model deployments, and digital twin templates, should only be released for clinical use after all necessary approvals for the clinical trial relevant to that investigator site have been received and after validation has been completed.

\subsubsection{4.3.6 System Failure}

Contingency procedures should be in place to prevent loss of data or loss of accessibility to data essential to participant safety, trial decisions, or trial outcomes in the event of physical AI system failure. These procedures should address:

\begin{description}[leftmargin=2.5cm,style=sameline]

\item[(a)] Robotic system failure during a clinical procedure, including safe shutdown sequences, data preservation, and clinical fallback procedures.

\item[(b)] AI inference engine failure, including reversion to human-only decision making and documentation of the period during which AI assistance was unavailable.

\item[(c)] Digital twin platform failure, including procedures for continuing clinical care without twin-based guidance and retroactive twin synchronisation when service is restored.

\item[(d)] Network failure affecting federated learning or multi-site data exchange, including local data buffering and deferred synchronisation.

\end{description}

\subsubsection{4.3.7 Technical Support}

\begin{description}[leftmargin=2.5cm,style=sameline]

\item[(a)] Mechanisms should be in place to document, evaluate, and manage issues with physical AI systems raised by clinical users, including robotic system operators, AI tool users, and digital twin managers. Periodic review of cumulative issues should identify repeated and systemic problems.

\item[(b)] Physical AI system defects and issues should be resolved according to their clinical criticality, with safety-critical issues requiring immediate resolution and documented root cause analysis.

\end{description}

\subsubsection{4.3.8 User Management}

\begin{description}[leftmargin=2.5cm,style=sameline]

\item[(a)] Access controls for physical AI systems used in clinical trials should limit system access to authorised users and ensure attributability to an individual. For robotic systems, this includes operator authentication before each clinical use session. For AI systems, this includes role-based access to model configuration and inference parameters. Security measures should achieve the intended level of protection without impeding clinical workflow.

\item[(b)] Procedures should ensure that user access permissions for physical AI systems are appropriately assigned based on the user's documented qualifications, clinical duties, and blinding arrangements. Access permissions should be revoked when they are no longer needed, including when personnel rotate off the trial or when specific physical AI systems are decommissioned.

\item[(c)] Authorised users and access permissions for all physical AI systems should be clearly documented, maintained, and retained. Records should include time-stamped grant and revocation of permissions.

\end{description}

\section{APPENDICES}

\section{Appendix A. PHYSICAL AI SYSTEM DOCUMENTATION}

A.1 Introduction

The Physical AI System Documentation package is a compilation of the technical, safety, and clinical data on the physical AI systems that are relevant to their use in oncology clinical trials with human participants. This documentation serves a purpose analogous to the Investigator's Brochure described in the prior ICH E6(R3) regulation Appendix A, extended to address robotic systems, AI models, digital twins, and simulation environments. Its purpose is to provide investigators and other trial personnel with the information needed to understand the capabilities, limitations, safety profiles, and proper use of each physical AI system specified in the protocol.

\paragraph{A.1.1} Development of the Physical AI System Documentation

The sponsor is responsible for ensuring that comprehensive Physical AI System Documentation is developed and maintained for each system used in the trial. This documentation should be reviewed at least annually and revised as significant new information becomes available, including software updates, hardware revisions, AI model retraining results, safety incident analyses, and new clinical experience data. The USL assessment (v1.4.0 through v1.8.0) should be updated when platform capabilities change materially.

\paragraph{A.1.2} Risk-Benefit Assessment for Physical AI Systems

The Physical AI System Documentation should include a balanced risk-benefit assessment for each system, covering mechanical risks, AI prediction risks, cybersecurity risks, and the anticipated clinical benefits. This assessment should be updated as new safety data and clinical evidence accumulate during the trial.

A.2 General Considerations

These considerations delineate the minimum information that should be included in the Physical AI System Documentation. The type and extent of information available will vary with the maturity of the physical AI system, its USL score, and the stage of clinical development.

The Physical AI System Documentation should include:

\paragraph{A.2.1} Title Page

This should provide the sponsor's name, the identity of each physical AI system (manufacturer, model designation, software version), the USL score with version reference, and the release date of the documentation. An edition number and the date of data inclusion cut-off should be provided.

\paragraph{A.2.2} Confidentiality Statement

The sponsor may include a statement instructing the investigator and other recipients to treat the Physical AI System Documentation as confidential, for the sole information and use of authorised trial personnel, regulatory authorities, and the IRB/IEC.

A.3 Contents of the Physical AI System Documentation

The documentation should contain the following sections:

\paragraph{A.3.1} Table of Contents

\paragraph{A.3.2} Summary

A brief summary (preferably not exceeding three pages) should highlight the significant technical, safety, and clinical information for each physical AI system, including its intended clinical use, key performance specifications, USL score summary, and known limitations.

\paragraph{A.3.3} System Description and Specifications

For each robotic system, the documentation should include:

\begin{description}[leftmargin=2.5cm,style=sameline]

\item[(a)] Mechanical specifications: degrees of freedom, workspace volume, payload capacity, positional repeatability, maximum speed, and weight.

\item[(b)] Control specifications: control frequency (Hz), force/torque sensing capabilities, collision detection parameters, and safety-rated monitored stop response time.

\item[(c)] Software specifications: operating system, control software version, communication protocols (e.g., ROS 2 Jazzy, EtherCAT), and supported programming interfaces.

\item[(d)] USL dimension scores: simulation switching (Dimension 1), AI integration (Dimension 2), cross-robot sharing (Dimension 3), and clinical trial collaboration (Dimension 4), with rationale for each score.

\end{description}

For each AI model, the documentation should include:

\begin{description}[leftmargin=2.5cm,style=sameline]

\item[(a)] Model architecture: type (e.g., vision-language-action, diffusion policy, transformer, convolutional neural network), parameter count, input and output specifications.

\item[(b)] Training data: description of training datasets, including size, source, demographic composition, de-identification status, and any known biases or limitations.

\item[(c)] Validation performance: metrics appropriate to the clinical task (e.g., accuracy, sensitivity, specificity, AUC, calibration error), evaluated on independent test sets representative of the target oncology population.

\item[(d)] Intended use and limitations: clear statement of the clinical contexts in which the model is validated for use, and explicit documentation of conditions under which the model may produce unreliable results.

\end{description}

For each digital twin, the documentation should include:

\begin{description}[leftmargin=2.5cm,style=sameline]

\item[(a)] Mathematical model: governing equations (e.g., reaction-diffusion for glioblastoma invasion, logistic growth, Gompertz model, PK/PD models for chemotherapy response), numerical methods, and computational requirements.

\item[(b)] Calibration methodology: patient data inputs required (e.g., longitudinal imaging, biomarker measurements), calibration algorithm, convergence criteria, and expected calibration time.

\item[(c)] Validation evidence: comparison of twin predictions against observed clinical outcomes in representative patient cohorts, including prediction accuracy for tumor growth, treatment response, and surgical outcomes.

\item[(d)] Clinical integration: data flow from clinical systems to the twin platform, prediction delivery to clinical decision makers, and human override mechanisms.

\end{description}

\paragraph{A.3.4} Safety Studies and Testing

The documentation should provide a summary of all safety studies conducted for each physical AI system, including:

\begin{description}[leftmargin=2.5cm,style=sameline]

\item[(a)] Mechanical safety testing for robotic systems, including collision force measurements, emergency stop response testing, workspace boundary enforcement verification, and compliance with ISO/TS 15066 (collaborative workspace safety) and IEC 80601-2-77 (robotically assisted surgical equipment).

\item[(b)] Software safety testing, including fault injection testing, graceful degradation verification, and compliance with IEC 62304 (medical device software lifecycle).

\item[(c)] AI model safety testing, including adversarial robustness evaluation, out-of-distribution detection, and failure mode analysis for clinically relevant edge cases.

\item[(d)] Cybersecurity testing, including penetration testing, vulnerability assessment, and compliance with FDA premarket cybersecurity guidance.

\item[(e)] Simulation validation results, including cross-framework behavioral comparison using the unification bridge (v1.0.0), with documented trajectory deviations, force discrepancies, and task completion rate differences across NVIDIA Isaac Lab, MuJoCo, Gazebo, and PyBullet environments.

\end{description}

\paragraph{A.3.5} Clinical Experience

The documentation should include a summary of all clinical experience with each physical AI system, including:

\begin{description}[leftmargin=2.5cm,style=sameline]

\item[(a)] Prior clinical trial experience: summary of completed and ongoing trials involving the system, including participant numbers, procedures performed, adverse events, and clinical outcomes.

\item[(b)] Non-trial clinical use: for systems with prior clinical deployment (e.g., da Vinci surgical system with 14 million or more procedures in 69 countries; 9,000 or more installed systems), a summary of real-world performance data, safety reports, and any recalls or field safety actions.

\item[(c)] Relevant research experience: summary of peer-reviewed research publications using the system, including the size and quality of the research evidence base (e.g., dVRK with 500 or more research publications across 45 or more institutions).

\end{description}

\paragraph{A.3.6} Summary of Data and Guidance

This section should provide an overall assessment of each physical AI system's readiness for clinical trial deployment, integrating mechanical safety data, AI model validation evidence, digital twin accuracy assessments, simulation validation results, and clinical experience. The assessment should include specific recommendations for the clinical investigator regarding system operation, monitoring, and safety precautions.

\section{Appendix B. CLINICAL TRIAL PROTOCOL FOR PHYSICAL AI TRIALS}

Physical AI oncology trials should be described in a clear, concise, and operationally feasible protocol that incorporates all specifications for physical AI systems in addition to standard clinical trial elements. The protocol should be designed to minimise unnecessary complexity and to mitigate or eliminate important risks to the rights, safety, and well-being of trial participants and the reliability of data. This appendix adapts the prior ICH E6(R3) Appendix B protocol template to address physical AI trial requirements.

The contents of a physical AI trial protocol should generally include the following topics:

B.1 General Information

B.1.1 Protocol title, unique protocol identifying number, and date. Any amendment(s) should also bear the amendment number(s) and date(s).

B.1.2 Name and address of the sponsor.

B.1.3 Name and title of the person(s) authorised to sign the protocol for the sponsor.

B.1.4 Physical AI System Inventory: a complete list of all physical AI systems specified in the protocol, including manufacturer, model, software version, and USL score for each system.

B.2 Background Information

B.2.1 Name and description of the investigational product(s) and the physical AI systems to be used.

B.2.2 A summary of findings from preclinical testing of physical AI systems, including simulation validation results, safety testing outcomes, and AI model performance data.

B.2.3 Summary of the known and potential risks and benefits of physical AI system interactions for trial participants, including risks from the prior ICH E6(R3) protocol requirements and additional physical AI-specific risks.

B.2.4 Description of the physical AI system configurations, operating parameters, and safety thresholds to be used in the trial.

B.2.5 A statement that the trial will be conducted in compliance with the protocol, Good Clinical Practice (GCP), the applicable regulatory requirements, and all physical AI system safety standards.

B.2.6 Description of the patient population, including any specific considerations for physical AI interactions (e.g., contraindications for robotic procedures, patient factors affecting digital twin calibration accuracy).

B.3 Trial Objectives and Purpose

A clear description of the scientific objectives and the purpose of the trial, including any objectives related to physical AI system performance evaluation, AI model validation, or digital twin methodology assessment.

B.4 Trial Design

B.4.1 Primary and secondary endpoints, including any physical AI system performance endpoints.

B.4.2 Description of the trial design, including any physical AI-specific design elements such as robotic system randomisation, AI model comparison arms, or digital twin-guided adaptive dosing.

B.4.3 Measures to minimise bias, including blinding of physical AI system configurations where appropriate.

B.4.4 Physical AI system specifications: detailed description of each system's operating parameters, safety thresholds, and acceptable performance ranges for the trial.

B.4.5 Schedule of events, including physical AI system calibration verification, AI model performance checks, and digital twin synchronisation points.

B.4.6 Expected duration of participant involvement, including time for physical AI system interactions.

B.5 Selection of Participants

B.5.1 Inclusion criteria, including any physical AI-specific eligibility requirements.

B.5.2 Exclusion criteria, including contraindications for physical AI system interactions.

B.5.3 Screening procedures, including any physical AI-related assessments (e.g., patient suitability for robotic procedure, digital twin calibration feasibility).

B.6 Discontinuation Criteria

B.6.1 When and how to discontinue participants from specific physical AI system interactions.

B.6.2 Physical AI system-specific stopping criteria (e.g., robotic safety limit breach, AI model confidence below threshold, digital twin prediction failure).

B.6.3 Procedures for continuing clinical care when physical AI components are discontinued.

B.7 Physical AI System Management

B.7.1 Detailed operating procedures for each physical AI system, including startup, calibration, clinical operation, and shutdown.

B.7.2 Acceptable and prohibited concurrent technologies or interventions during physical AI system operation.

B.7.3 Strategies to monitor participant adherence to physical AI-related instructions and protocols.

B.8 Assessment of Efficacy

B.8.1 Specification of efficacy parameters, including any physical AI system performance metrics that serve as efficacy endpoints.

B.8.2 Methods and timing for assessing efficacy, including integration of physical AI data with clinical outcome assessments.

B.9 Assessment of Safety

B.9.1 Specification of safety parameters, including physical AI system-specific safety metrics (e.g., robotic contact forces, AI prediction confidence, digital twin accuracy).

B.9.2 Methods for recording and assessing physical AI safety events.

B.9.3 Procedures for reporting physical AI system adverse events.

B.9.4 Follow-up procedures for participants who experienced physical AI system-related adverse events.

B.10 Statistical Considerations

B.10.1 Statistical methods, including any methods specific to analysing physical AI system data (e.g., mixed-effects models accounting for robotic system variability, Bayesian updating of digital twin predictions).

B.10.2 Sample size justification, including consideration of physical AI system variability.

B.10.3 Analysis of physical AI covariates: pre-specified analyses examining the influence of robotic system type, AI model version, and digital twin calibration quality on primary and secondary endpoints.

B.11 Direct Access to Source Records

The protocol should specify that investigators provide direct access to source records for physical AI systems, including robotic telemetry, AI inference logs, and digital twin state data, for the purposes of monitoring, auditing, and regulatory inspection.

B.12 Quality Control and Quality Assurance

B.12.1 Description of critical to quality factors for physical AI system components.

B.12.2 Monitoring approaches for physical AI systems.

B.12.3 Noncompliance handling procedures for physical AI system specifications.

B.13 Ethics

Description of ethical considerations relating to the trial, including ethics of AI-assisted clinical decisions, participant autonomy in robotic procedures, and equitable access to physical AI-enabled care.

B.14 Data Handling and Record Keeping for Physical AI Data

B.14.1 Specification of physical AI data to be collected, including data types, sampling rates, and storage formats.

B.14.2 Identification of physical AI data considered to be source records.

B.14.3 Retention requirements for physical AI system records.

B.15 Financing and Insurance

Financing, insurance, and indemnification arrangements, including coverage for physical AI system-related injuries.

B.16 Publication Policy

Publication policy, including provisions for reporting physical AI system performance data.

\section{Appendix C. ESSENTIAL RECORDS FOR PHYSICAL AI CLINICAL TRIALS}

C.1 Introduction

C.1.1 Physical AI oncology clinical trials generate essential records beyond those of traditional clinical trials, including robotic system documentation, AI model artifacts, digital twin models, simulation validation records, and federated learning logs. The nature and extent of these records are dependent on the trial design, the physical AI systems used, and the risk-proportionate approach applied.

C.1.2 Determining which physical AI records are essential should follow the guidance in this appendix, supplementing the criteria from the prior ICH E6(R3) Appendix C.

C.1.3 Essential physical AI records permit and contribute to the evaluation of the conduct of the trial in relation to the compliance of the investigator and sponsor with Good Clinical Practice, applicable regulatory requirements, and physical AI system safety standards. These records are used for monitoring, auditing, and regulatory inspection purposes.

C.2 Management of Essential Physical AI Records

C.2.1 Physical AI records should be identifiable, version controlled, and should include authors, reviewers, and approvers along with date and signature where necessary.

C.2.2 For physical AI system activities delegated to service providers (e.g., robotic system maintenance, AI model training, digital twin development), arrangements should be made for access and management of essential records throughout the trial and for retention following trial completion.

C.2.3 Essential physical AI records should be maintained in or referred to from the trial master file (TMF) and investigator site file (ISF), supplemented by physical AI system-specific repositories for robotic telemetry, AI model artifacts, and digital twin data.

C.2.4 The storage systems used for physical AI records should provide for appropriate identification, version history, search, and retrieval, with attention to the large data volumes generated by robotic telemetry and sensor measurements.

C.2.5 Physical AI records should be collected and filed in a timely manner. Some records (e.g., robotic system validation, AI model documentation, simulation validation reports) should be in place prior to the start of the trial.

C.2.6 Physical AI records should be retained in a way that ensures they remain complete, readable, and readily available for the required retention period. Alteration should be traceable.

C.3 Essentiality of Physical AI Trial Records

C.3.1 In addition to the criteria from the prior ICH E6(R3) Appendix C, a physical AI trial record is essential if it:

\begin{description}[leftmargin=2.5cm,style=sameline]

\item[(a)] Documents the identity, configuration, calibration, and validation status of each physical AI system used in the trial.

\item[(b)] Documents the version, training provenance, validation performance, and clinical deployment history of each AI model used in the trial.

\item[(c)] Documents the mathematical basis, calibration methodology, validation evidence, and prediction record of each digital twin model used in the trial.

\item[(d)] Documents simulation validation results, including cross-framework behavioral comparisons and identified discrepancies.

\item[(e)] Documents physical AI system safety events, including robotic incidents, AI prediction failures, digital twin accuracy deviations, and cybersecurity incidents.

\item[(f)] Documents training and competency assessments for all personnel operating or overseeing physical AI systems.

\item[(g)] Documents federated learning participation, including aggregation rounds, privacy budget expenditure, and data harmonization actions.

\item[(h)] Documents maintenance, calibration, and repair activities for robotic systems throughout the trial.

\item[(i)] Documents the end-of-trial disposition of physical AI systems, including data extraction, secure erasure of participant data, and system return or decommissioning.

\end{description}

C.3.2 The following Physical AI Essential Records Table provides a non-exhaustive list of records that should be retained when produced:

Physical AI Essential Records Table

Note: An asterisk (*) identifies records that should generally be in place prior to the start of the trial.

\begin{description}[leftmargin=0.5cm,style=sameline]

\item[--] Physical AI System Documentation package for each system*
\item[--] USL assessment report for each robotic platform*
\item[--] Robotic system validation and calibration records*
\item[--] AI model validation reports with performance metrics*
\item[--] Digital twin validation reports with prediction accuracy data*
\item[--] Simulation validation reports (cross-framework comparison)*
\item[--] Cybersecurity assessment and penetration testing reports*
\item[--] Physical AI system operator training and competency records*
\item[--] Robotic system maintenance and calibration logs during trial
\item[--] AI model version control logs and update documentation
\item[--] Digital twin calibration and recalibration records
\item[--] Robotic system telemetry and procedure logs
\item[--] AI model inference logs and decision records
\item[--] Digital twin state histories and prediction records
\item[--] Physical AI safety event reports and root cause analyses
\item[--] Cybersecurity incident logs and response records
\item[--] Federated learning round logs, privacy budget records, and aggregation documentation
\item[--] Physical AI monitoring reports (site and centralised)
\item[--] Physical AI system decommissioning and data erasure records
\item[--] Signed agreements with physical AI system service providers*

\end{description}

\section{GLOSSARY}

The following definitions supplement the glossary from the prior ICH E6(R3) regulation with terms specific to physical AI oncology clinical trials. Standard clinical trial terms (Adverse Event, Audit, Blinding, Clinical Trial, Informed Consent, Investigator, Sponsor, etc.) retain their definitions from the prior ICH E6(R3) regulation and are not repeated here.

Agentic AI: Artificial intelligence systems capable of autonomous reasoning, planning, and action execution with minimal human prompting. In oncology clinical trials, agentic AI includes LLM-based surgical assistants, multi-agent coordination systems (e.g., CrewAI v1.6.1, LangGraph v1.1.0), and autonomous simulation orchestrators. Agentic AI systems operate under human oversight as specified in the human oversight quality management framework (v1.2.0).

AI Model Drift: The degradation of AI model performance over time due to changes in the input data distribution, population characteristics, or clinical context. Detection and management of drift is a critical monitoring activity in physical AI trials.

Calibration (Robotic): The process of measuring and adjusting a robotic system's sensors, actuators, and coordinate systems to ensure that the system's internal model accurately represents its physical configuration. Calibration verification should be performed before each clinical use session.

Calibration (AI Model): The degree to which an AI model's predicted probabilities match the observed frequencies of outcomes. A well-calibrated model produces predictions that are neither overconfident nor underconfident.

Cobot (Collaborative Robot): A robotic system designed to operate alongside human workers in a shared workspace, equipped with safety features including force limiting, speed reduction, and collision detection. In oncology trials, cobots assist with sample handling, drug preparation, and laboratory automation. Evaluated cobots include Franka Emika Panda (USL 7.4), Kinova Gen3 (USL 5.7), and UFACTORY xArm 7 (USL 3.4).

Cross-Framework Validation: The process of verifying that a robotic policy or behavior trained in one simulation framework produces equivalent results when executed in a different simulation framework. The unification bridge (v1.0.0) provides tools for cross-framework validation between NVIDIA Isaac Lab, MuJoCo, Gazebo, and PyBullet.

Cybersecurity Incident: An event that compromises the confidentiality, integrity, or availability of a physical AI system or its data. In the context of clinical trials, cybersecurity incidents include unauthorized access to robotic controllers, tampering with AI model parameters, corruption of digital twin data, and network-based attacks on clinical infrastructure.

Differential Privacy: A mathematical framework for quantifying and limiting the privacy loss when publishing statistical analyses of sensitive data. In federated oncology trials, differential privacy mechanisms (epsilon and delta parameters) are applied to gradient updates and aggregate statistics to protect individual patient data. The federation framework (v1.0.0) implements Gaussian and Laplacian mechanisms for differential privacy.

Digital Twin: A computational model of a physical entity (patient, tumor, organ, or treatment process) that is continuously updated with real-world data to provide predictions, simulations, and decision support. In oncology trials, digital twins model patient-specific tumor growth, treatment response, and surgical anatomy. The digital twins framework (v1.0.0 through v1.3.0) provides tools for creating, calibrating, and validating patient-specific digital twins.

Diffusion Policy: A generative AI technique that uses denoising diffusion probabilistic models to generate robotic trajectories, producing smooth, multi-modal action distributions for complex manipulation tasks such as tumor resection and needle insertion.

Domain Randomization: A simulation training technique that varies environmental parameters (lighting, textures, physics properties, sensor noise) during training to produce robotic policies that transfer more reliably from simulation to real-world deployment.

Edge Computing: Computation performed on devices located close to the data source (e.g., at the robotic system or in the operating room) rather than in a centralised data centre. Edge computing enables low-latency AI inference for real-time robotic control and is relevant to the deployment architecture of physical AI clinical trials.

Federated Learning: A machine learning approach that trains models across multiple sites without centralising raw data. Each site trains on its local data and shares only model updates (gradients or weights) with a central aggregator. The federation framework (v1.0.0) supports FedAvg, FedProx, and SCAFFOLD algorithms for multi-site oncology trials.

Force Limiting: A safety mechanism in robotic systems that restricts the maximum force the robot can exert on a person or object. Force limits are a critical safety parameter in clinical robotic procedures and must be validated and monitored throughout the trial.

GR00T N1.6: NVIDIA's humanoid robot foundation model providing vision-language-action capabilities for general-purpose robotic control. GR00T N1.6 is referenced in the physical-ai-oncology-trials repository (v2.2.0) for humanoid robot integration.

Humanoid Robot: A full-body robotic system with a form factor resembling the human body, typically including a torso, two arms, two legs, and a head. In oncology trials, humanoid robots may serve in patient assistance, logistics, and pediatric companionship roles. Evaluated humanoid robots include Boston Dynamics Atlas (USL 5.8), Agility Digit (USL 4.2), and Tesla Optimus (USL 3.6).

Model Context Protocol (MCP): A standardized protocol for communication between AI agents and external tools, governed by the Linux Foundation AI and Data Foundation (AAIF). In physical AI oncology trials, MCP enables standardized integration between AI assistants and robotic systems, clinical databases, and simulation environments.

Physical AI: The integration of artificial intelligence with physical robotic systems and sensor technologies to perform tasks in the real world. Physical AI encompasses the full pipeline from simulation-based training through real-world deployment of intelligent robotic systems.

Physical AI System Documentation: The technical, safety, and clinical documentation package for a physical AI system, serving a purpose analogous to the Investigator's Brochure from the prior ICH E6(R3) regulation. See Appendix A.

Reinforcement Learning (RL): A machine learning paradigm in which an agent learns to make decisions by taking actions in an environment and receiving rewards or penalties. In oncology clinical trials, RL is used for surgical autonomy training, treatment dose optimization, and adaptive protocol design. GPU-accelerated RL training via NVIDIA Isaac Lab (v2.3.1) enables training in parallel environments for surgical tasks.

Secure Aggregation: A cryptographic protocol that allows multiple parties to compute the sum of their private inputs without revealing individual values. In federated oncology trials, secure aggregation protects site-specific model updates during the aggregation step. The federation framework (v1.0.0) implements additive secret sharing and pairwise masking protocols.

Sim-to-Real Transfer: The process of deploying a robotic policy trained in simulation to a physical robot in the real world. The sim-to-real gap refers to the performance difference between simulation and real-world execution, which must be minimised and documented for clinical applications.

Simulation Framework: A software platform for simulating robotic systems, physics interactions, and sensor observations. Recognized simulation frameworks include NVIDIA Isaac Lab (v2.3.1), Isaac Sim (v5.0.0), MuJoCo (v3.4.0), MuJoCo Warp (Beta), Gazebo Sim (Jetty, v10.0.0), and PyBullet (v3.2.5).

Surgical Robot: A robotic system designed for use in surgical procedures, typically operated via teleoperation with surgical console controls. In oncology trials, surgical robots perform tumor resection, biopsy, and minimally invasive procedures. Evaluated surgical robots include da Vinci (dVRK, USL 7.1), Hugo RAS (USL 4.5), and Versius (USL 3.4).

Unification Standard Level (USL): A scoring framework (v1.4.0 through v1.8.0) for evaluating robotic platform readiness for unified, multi-site oncology clinical trials. USL scores range from 1.0 to 10.0 across four dimensions: simulation switching, AI integration, cross-robot sharing, and clinical trial collaboration. Scores map to ten levels from Minimal (1.0) to Reference (10.0) and three bands: Foundational (1.0 to 4.9), Intermediate (5.0 to 6.9), and Advanced (7.0 to 10.0).

Vision-Language-Action (VLA) Model: A multimodal AI model that takes visual observations and natural language instructions as input and produces robotic actions as output. VLA models such as NVIDIA GR00T N1.6 enable natural language programming of robotic behaviors in clinical settings.

\clearpage
\nocite{ich-e6-r3}
\printbibliography[title={Bibliography}]
\end{document}
